\section{Human-AI Collaboration Process}
\label{sec:collaboration}

The framework was built through intensive interactive Claude Code sessions in which the human drove scientific strategy and acceptance criteria, and the AI drove implementation, parallelization, and documentation. The collaboration is bidirectional: the human curates the corpora, defines what counts as a passing user-acceptance test, and corrects framing drifts; the AI implements, tests in parallel across multiple scenarios, and surfaces inconsistencies for review. This section makes three observations specific to how that collaboration shaped the current framework.

\subsection{The GSD workflow as methodology}

The framework was built over twenty phases under a structured planning workflow (``Get Shit Done,'' GSD), in which each phase begins with a written goal and acceptance criteria, proceeds through a discussion-and-research stage, lands a plan, executes against atomic commits, and verifies goal achievement before closing. The workflow's artifacts (\texttt{ROADMAP.md}, \texttt{REQUIREMENTS.md}, per-phase \texttt{PLAN.md} and \texttt{VERIFICATION.md}) are local-only by design (gitignored under \texttt{.planning/}) but they shape the public commit history: each commit message names its phase and the requirement it discharges, producing an auditable trail from goal to implementation to verification.

The twenty phases that produced the v3.0 framework grouped into three arcs. Phases~1--10 were foundational: domain configuration, document ingestion, entity extraction and graph construction, cross-reference analysis, the interactive dashboard, repo reorganization, the testing framework, the domain-creation wizard, consumer decoupling, and documentation refresh. Phase~11 was end-to-end scenario validation, the V2 rebuilds of S1--S5 reported in Table~\ref{tab:scenarios}. Phases~12--20 produced the v3.0 graph-fidelity-and-honest-limits arc: structural-biology document signatures (Phase~12), extraction pipeline reliability and Pydantic-strict write-time validation (Phase~13, FIDL-02c), sentence-aware chunk overlap via \texttt{chonkie} (Phase~14, FIDL-03; the chunking strategy described in Section~\ref{sec:architecture}), format discovery parity (Phase~15, FIDL-04), wizard sample-window expansion (Phase~16, FIDL-05), domain awareness in downstream consumers (Phase~17, FIDL-06), per-domain epistemic and validator extensibility (Phase~18, FIDL-07; the Layer 2 customization paths described in Section~\ref{sec:epistemic}), wizard and CLI ergonomics (Phase~19, FIDL-08), and the public-facing capacity-and-limits documentation (Phase~20, FIDL-09). Each requirement was scoped tightly enough to verify in isolation; their composition produced the v3.0 release. Phase~21 added the clinicaltrials domain (v3.1), and Phase~22 added the fda-product-labels domain (v3.2).

\subsection{External contribution: the controlled-port pattern}

The clinicaltrials domain (Section~\ref{sec:domain-clinicaltrials}) and the fda-product-labels domain (Section~\ref{sec:domains}) were initially contributed by Chris Davidson via pull request. In both cases the original contribution arrived against a stale fork (the contributor's main branch lagged the upstream main by a release or more), and a direct rebase or merge would either lose context or pull in unrelated drift. The framework adopted a \emph{controlled-port pattern}: rather than rebasing the contributor's branch onto upstream main, the maintainer creates a new branch off current upstream main, hand-ports the substantive contribution file-by-file with attribution preserved via Git's \texttt{Co-Authored-By} trailer, and reviews each ported file against current conventions. This avoided two recurring failure modes (contributor-introduced \texttt{.planning/} files leaking into the public repo, and stub configurations replacing hand-tailored files) and preserved authorship attribution. The pattern is documented in repository memory for future contributions.

\subsection{Honesty as a development practice}

Several v3 commits are corrections rather than additions. The README was found to overclaim ``domain knowledge compounds across scenarios -- patterns generalize, edge cases surface, the analyst persona learns what to flag,'' implying an automatic-learning mechanism that the framework does not implement. The correction reframed scenario count as a track record (pressure-testing depth) rather than an auto-learning signal, and opened GitHub Issue~\#15 to track the aspirational mechanism. The visible commit (\texttt{a8a91db}) does not add a feature; it walks back a marketing claim. The collaboration model includes the human's role in catching such drifts -- the AI alone is unlikely to identify when its own framing has shifted from ``claim'' to ``promise.''

The pattern recurs in three places: the README compounding correction, the demo-script accuracy checklist (every claim mapped to a verifiable codebase source), and this paper's Section~\ref{sec:evaluation} (``What the evaluation does not establish''). Each is an artifact of treating documentation honesty as a first-class development task rather than as polish.
